\chapter{Paradoxos de Representação e Formalização}
\label{chap:apendiceRepresentacaoFormalizacao}

Adicionar Introdução

\section{O Paradoxo de Skolem (1922)}

Apresentar o Teorema de Löwenheim-Skolem.

% Fontes:
%   - Skolem, artigo original traduzido em:
%     van Heijenoort, From Frege to Gödel, Harvard University Press, 1967
%   - Enderton, A Mathematical Introduction to Logic, Academic Press, 2001
%   - Cohen, Set Theory and the Continuum Hypothesis, Dover, 2008

\section{O Paradoxo de König (1905)}

Apresentar o episódio histórico no Congresso Internacional de
Matemática de 1904.

% Fontes:
%   - Hausdorff, Set Theory, Chelsea, 1962 (já na bibliografia do trabalho)
%   - Moore, Zermelo's Axiom of Choice, Springer, 1982

\section{O Teorema de Goodstein (1944)}

Apresentar a sequência de Goodstein e seu crescimento explosivo
seguido de colapso inevitável a zero. Demonstrar a prova via
ordinais transfinitos: associar a cada termo um ordinal em
notação hereditária com \(\omega\) no lugar da base e usar a
boa-ordenação dos ordinais para concluir. Conectar com o
Capítulo~7 (Números Ordinais). Apresentar o resultado de
Kirby e Paris (1982): a terminação não é demonstrável dentro
da aritmética de Peano,
situando o teorema na mesma família do resultado de Gödel
do Apêndice~B, mas com enunciado elementar e papel dos
ordinais completamente explícito.

% Fontes:
%   - Goodstein, R. L. On the restricted ordinal theorem.
%     Journal of Symbolic Logic, v. 9, n. 2, p. 33–41, 1944.
%   - Kirby, L.; Paris, J. Accessible independence results for
%     Peano arithmetic. Bulletin of the London Mathematical
%     Society, v. 14, n. 4, p. 285–293, 1982.
%   - Cichon, E. A. A short proof of two recently discovered
%     independence results using recursion theoretic methods.
%     Proceedings of the American Mathematical Society,
%     v. 87, n. 4, p. 704–706, 1983.